\documentclass{article}

\usepackage[utf8]{inputenc}
\usepackage[T1]{fontenc}
\usepackage{amsmath}
\usepackage{graphicx}
\usepackage{caption}
\usepackage{float}
\usepackage[top=1.5cm, bottom=1.5cm]{geometry}
\setlength{\parskip}{1em}
\setlength{\parindent}{0pt}
\captionsetup{justification=centering}

\title{Sortowanie plików sekwencyjnych metodą dużych buforów}
\author{Igor Sitko-Bajorski}
\date{Listopad 2025}

\begin{document}

\maketitle

\section{Opis metody}
Sortowanie plików sekwencyjnych metodą dużych buforów zakłada ograniczoną ilość pamięci operacyjnej w stosunku do rozmiaru pliku. W celu dostępu do danych w pamięci RAM wykorzystujemy zdefiniowaną ilość buforów, \textbf{oznaczoną \textit{n}}, o określonej pojemności (w tym przypadku, określonym rozmiarem strony dyskowej). Ilość rekordów mieszczących się w jednym buforze nazywamy czynnikiem blokującym, oraz \textbf{oznaczamy \textit{b}}, gdzie:
\[
b = \text{(Rozmiar Strony Dyskowej ) / (Rozmiar Rekordu)}
\]

Sortowanie można rozdzielić na dwa kroki:
\begin{enumerate}
    \item Etap tworzenia serii
    \item Etap złączenia serii
\end{enumerate}

Etap tworzenia serii polega na:
\begin{enumerate}
    \item wczytaniu do buforów zawartości pliku,
    \item posortowaniu \(n \cdot b\) rekordów w buforach, algorytmem sortowania wydajnym pamięciowo (w tym przypadku QuickSort),
    \item nadpisania wczytanych danych posortowanymi. Te czynności są powtarzane do momentu osiągnięcia końca pliku.
\end{enumerate}

Etap złączenia serii polega na: 
\begin{enumerate}
    \item stworzeniu tymczasowego pliku wyjściowego,
    \item rozdzieleniu buforów: \textit{n - 1} pierwszych buforów przeznaczamy na dane wejściowe, ostatni na zapis wyników złączenia,
    \item wczytaniu do buforów początków posortowanych serii,
    \item utworzenie minimalnego kopca z pierwszych elementów w buforach,
    \item ekstrakcję minimum, i jego zapis do bufora wyjściowego,
    \item dodanie do kopca kolejnej wartości z bufora który zawierał poprzednie minimum.
\end{enumerate}

W przypadku wyczerpania danych w buforze wejściowym - wczytujemy do niego kolejną stronę serii.

W przypadku osiągnięcia końca serii bufora wejściowego nie dodajemy kolejnej wartości do kopca.

W przypadku zapełnienia danych w buforze wyjściowym - zapisujemy jego zawartość do pliku wyjściowego 

Kroki 5 i 6 są powtarzane do momentu opróżnienia kopca.

Te czynności są powtarzane do momentu osiągnięcia końca pliku. Pomiędzy wywołaniami zamieniamy rolami plik wejściowy i wyjściowy.

\section{Specyfikacja formatu pliku}
\subsection{Specyfikacja rekordu}
Pojedynczy rekord, opisuje obiekt na podstawie dwóch wartości masy oraz prędkości. Rekord składał się więc z dwóch wartości zmiennoprzecinkowych w formacie 64bitowym - \textit{double}. Pojedynczy rekord opisuje obiekt na podstawie masy oraz prędkości, 
a jego wartość jest wyznaczana na podstawie energii kinetycznej:
\[
E_k = \frac{mv^2}{2}.
\]

\subsection{Specyfikacja pliku}
Sortowany plik składał się z zapisanych binarnie, upakowanych bez separatora, rekordów o stałej długości.
\[
mvmv...mv
\]

Na każdy rekord przypada wówczas 16B = 8B (masa) + 8B (prędkość). Podczas odczytu, mieszczące się na stronie pełne rekordy zostają wpisane do bufora wyjściowego jako surowe bajty, które w celu odczytania należy skonwertować na wartości \textit{double}.

\section{Prezentacja wyników działania programu}
W formie podstawowej program wyświetla dwie informacje o procesie sortowania, liczbę faz etapu sortowania i etapu złączenia, oraz ilość operacji dyskowych.

W formie jawnej (verbose) wyświetla dodatkowo zawartość pliku (wartości rekordów) przed i po posortowaniu, a także zawartość pliku po każdej z faz.

W trakcie sortowania po każdym przejściu pętli w funkcjach \textbf{create\_runs()} oraz \textbf{merge()} zwiększany jest licznik faz danej operacji. Klasa \textbf{FileHandler} posiada wewnętrzny licznik operacji odczytu i zapisu na dysk, a także metodę \textbf{print()}, pozwalającą na odczyt wartości rekordów w pliku.

\newpage
\section{Eksperyment}
\subsection{Sposób przeprowadzenia eksperymentu}
Przyjęto wartości stałych:
\[
\begin{aligned}
R &= 16\text{B}\\
b &= 32\\
n &= 10
\end{aligned}
\]

Program przed rozpoczęciem sortowania generował rozmiar pliku na podstawie wzoru:
\[
\text{N} = it \cdot 100{,}000 + \left( \text{rand()} \bmod 5{,}000 \right),
\]
gdzie \textit{it} określa numer iteracji eksperymentu, \(0 \le it \le 10\).


Po określeniu rozmiaru, generował plik wejściowy o zadanej wielkości i zapisywał rozmiar (N) do pliku opisującego eksperyment.
Następnie zapisywał czas startu wywołania funkcji \textbf{sort()} i uruchamiał proces sortowania. W tym kroku, program dopisywał do pliku ilość wykonanych w tym procesie operacji dyskowych, będących sumą operacji r/w na sortowanym pliku i na pliku tymczasowym etapu złączenia. Dopisywał również ilość faz poszczególnych etapów. Po tym zapisywał czas końca wywołania, i dopisywał czas trwania wywołania do pliku opisującego eksperyment.


\subsection{Wyniki eksperymentu}
\subsubsection{Ilość operacji r/w}
Etap sortowania zajmie \( \frac{N}{b} \) operacji odczytu danych z pliku i tyle samo operacji zapisu - musimy 1 raz wczytać i zapisać cały plik.

Każda iteracja etapu złączenia również zajmie \( \frac{N}{b} \) operacji odczytu i zapisu danych z/do pliku.

Po pierwszej iteracji rozmiar serii będzie wynosił \( n \cdot b \). Po drugiej będzie to już \( n^2 \cdot b \) - złączamy (n - 1) serii, ale n >> 1, więc można założyć że \( n \approx n-1 \). Z tego wynika że tych etapów będzie ok. 
\[ 
    \lceil log_n(\frac{N}{nb}) \rceil = \lceil log_n(\frac{N}{b}) \rceil - 1
\]

Liczbę operacji dyskowych w zależności od rozmiaru pliku można oszacować więc wzorem:
\[
    \frac{2N}{b} \cdot log_n(\frac{2N}{b}) = \frac{2N}{b \cdot log_2(n)} \cdot log_2(\frac{N}{b})
\]

Dokładny wzór:
\[
    \lceil \frac{2N}{b} \rceil \cdot (1 +  \lceil log_{n - 1}(\frac{N}{nb}) \rceil)
\]

\newpage
Różnice między wartością aproksymowaną 1 wzorem, a rzeczywistą pokazuje poniższy wykres.

\begin{figure}[H]
    \centering
    \includegraphics[width=0.9\linewidth]{rw.png}
    \caption{Rzeczywiste i aproksymowane ilości operacji dyskowych. 
    Na osi \textit{X} rozmiar pliku~$N$, na osi \textit{Y} liczba operacji dyskowych}
    \label{fig:sort_operations}
\end{figure}

Na wykresie widać że zgodnie z przewidywaniem wykres rośnie prawie liniowo. Zaobserwowane ilości operacji dyskowych dla mniejszych rozmiarów pokrywają się, natomiast zwiększając rozmiar pliku, od ok. 100.000 rekordów występują rozbieżności.

Różnice między wartością aproksymowaną 2 wzorem, a rzeczywistą pokazuje poniższy wykres.

\begin{figure}[H]
    \centering
    \includegraphics[width=0.9\linewidth]{rw2.png}
    \caption{Rzeczywiste i aproksymowane ilości operacji dyskowych. 
    Na osi \textit{X} rozmiar pliku~$N$, na osi \textit{Y} liczba operacji dyskowych}
    \label{fig:sort_operations}
\end{figure}

Po dodaniu zaokrągleń i dokładnego logarytmu o podst. n - 1 do drugiej części, wzór na aproksymację dokładnie pokrywa się z obserwacjami.

\subsubsection{Ilość faz}
Ilości faz sortowania, to kolejna metryka którą przedstawia poniższy wykres.

\begin{figure}[H]
    \centering

    \begin{minipage}{0.48\linewidth}
        \centering
        \includegraphics[width=\linewidth]{s_phases.png}
        \caption{Fazy sortowania w zależności od rozmiaru pliku}
        \label{fig:sort_phases}
    \end{minipage}
    \hfill
    \begin{minipage}{0.48\linewidth}
        \centering
        \includegraphics[width=\linewidth]{m_phases.png}
        \caption{Fazy złączenia w zależności od rozmiaru pliku}
        \label{fig:merge_phases}
    \end{minipage}
\end{figure}

Ilość faz sortowania zwiększa się liniowo wraz z pojemnością pliku co zgadza się z teorią wzoru.

Ilość faz złączenia zgadza się z wartością teorytyczną.

\subsubsection{Czas wykonania}
\begin{figure}[H]
    \centering
    \includegraphics[width=0.9\linewidth]{time.png}
    \caption{Rzeczywiste i aproksymowane ilości operacji dyskowych. 
    Na osi \textit{X} rozmiar pliku na osi \textit{Y} czas wykonania}
    \label{fig:exec_time}
\end{figure}

Czas wykonania rośnie liniowo w zależności od rozmiaru pliku.

\subsection{Wnioski}
\subsubsection{Ilość operacji r/w}
Wzór bez zaokrągleń generuje większe rozbieżności dla operacji r/w, ponieważ założenie n >> 1 nie jest spełnione. Dla n = 10, 1 to już 10\% odchylenie. Ponadto błędy nawarstwiają się w przypadku zmiany podstawy logarytmu na dwójkowy.
Wykres posiada również załamanie spowodowane przeskokiem z ok. 5,000 próbek na ok. 100,000.

Stosując za to dokładny wzór jesteśmy w stanie, za każdym razem dokładnie określić ilość operacji dyskowych zgodną z rzeczywistością, w każdym punkcie wynik jest zgodny.

\subsubsection{Ilość faz}
Ilość faz sortowania wzrasta liniowo z podaną zależnością. Ilość faz jest równa \( \frac{N}{nb} \), bo w każdym kroku odczytamy \(nb\) danych (pojemność jednego bufora razy ilość buforów).

Ilość faz złączenia wzrasta logarytmicznie, zgodnie z podaną zależnością \( \lceil log_{n-1}(\frac{N}{nb}) \rceil \), a obserwacje zgadzają się ze wskazaniem teorytycznym.

\subsubsection{Czas wykonania}
Czas wykonania algorytmu zależy głównie od ilości operacji dyskowych, ponieważ to one są obarczone największym opóźnieniem. Pokrywa się to z danymi na wykresie, który wzrasta liniowo podobnie do ilości operacji odczytu i zapisu.

\end{document}
